\documentclass[11pt]{article}
\usepackage[a4paper,margin=0.65in]{geometry}
\usepackage[T1]{fontenc}
\usepackage{amsmath,amssymb,amsfonts}
\usepackage{graphicx}
\usepackage{pdfpages}
\usepackage[english,shorthands=off]{babel}
\usepackage{fancyhdr}
\usepackage[bookmarks=false,colorlinks=false]{hyperref}
\usepackage[protrusion=true,expansion=false]{microtype}
\emergencystretch=3em
\tolerance=600
\providecommand{\modd}{\mathrm{mod}}
\providecommand{\Disc}{\mathrm{Disc}}
\providecommand{\canont}{}
\providecommand{\asterism}{*}
\providecommand{\Hessian}{H}
\providecommand{\tome}{}
\providecommand{\page}{p.}
\pagestyle{fancy}\fancyhf{}\fancyhead[L]{Pages 426--437}\fancyhead[R]{Cayley --- Collected Papers, Vol. I}\renewcommand{\headrulewidth}{0.4pt}
\setlength{\headheight}{15pt}

\begin{document}

\begin{center}
{\LARGE\bfseries The Collected Mathematical Papers of}\\[0.5em]
{\LARGE\bfseries Arthur Cayley}\\[1em]
{\Large Volume I}\\[1em]
{\large PDF Pages 426--437 \quad (Book pages 404--415)}\\[2em]
\end{center}

\bigskip

% ===========================================================================
% PDF page 426 -- Book page 404 -- Paper [67] continued: NOTE SUR LES FONCTIONS ELLIPTIQUES
% ===========================================================================
\begin{center}
\textit{[Paper 67 concluded]}\\[0.5em]
\textbf{\large NOTE SUR LES FONCTIONS ELLIPTIQUES.}
\end{center}

\medskip

On aurait une valeur assez g\'en\'erale de $z$ en multipliant les diff\'erentes fonctions
$z_{i}$, chacune par une constante arbitraire, et en sommant les produits; mais dans le cas
actuel o\`u $z$ d\'enote le d\'enominateur de $\sqrt{\sin am\,nu}$, la valeur convenable de $z$ se
r\'eduit \`a
\[
z = z_{0} \pm z_{1} \cdots \pm z_{s} \cdots,
\]
o\`u $s$ est un nombre entier et positif, entre z\'ero et $\tfrac{1}{2}n$ ou $\tfrac{1}{2}(n - 1)$. On aura par
exemple dans le cas $n = 5$ (les signes \'etant tous positifs si $n$ est impair, et alternativement
positifs et n\'egatifs si $n$ est pair), en supprimant les puissances n\'egatives de $a$
(lesquelles s'entred\'etruisent):
\begin{align*}
z_{0} &= 1,\\
z_{1} &= 64a^{2}x^{17} - a^{2}(160x^{13} + 240x^{17}) + a(140x^{9} + 368x^{13} + 360x^{17})\\
&\qquad\qquad\qquad\qquad\qquad\qquad - (50x^{5} + 125x^{9} + 300x^{13} + 275x^{17}),\\
z_{2} &= 16a^{2}x^{4} - a(80x^{0} + 20x^{4}) + (170x^{0} + 62x^{4} + 5x^{8}),
\end{align*}
et de l\`a:
\begin{align*}
z &= 1 - 50x + 140ax^{5} - (125 + 160a^{2})x^{9} + (368a + 64a^{3})x^{13} - (300 + 240a^{2})x^{17}\\
&\quad + 360ax^{21} - 105x^{25} - 80ax^{29} + (62 + 16a^{2})x^{33} - 20ax^{37} + 5x^{41};
\end{align*}
ce qui est effectivement la valeur de $z$ que j'ai trouv\'ee dans le m\'emoire cit\'e pour
ce cas particulier.

\vfill\pagebreak

% ===========================================================================
% PDF page 427 -- Book page 405 -- Paper [68] begins
% ===========================================================================
\begin{center}
\textbf{\Large 68.}\\[1em]
\textbf{\large ON THE APPLICATION OF QUATERNIONS TO THE THEORY OF}\\[0.3em]
\textbf{\large ROTATION.}
\end{center}

\medskip

\begin{center}
\textit{[From the Philosophical Magazine, vol.\ xxxiii.\ (1848), pp.\ 196--200.]}
\end{center}

\medskip

In a paper published in the \textit{Philosophical Magazine}, February 1845, [20], I showed
how some formulae of M.\ Olinde Rodrigues relating to the rotation of a solid body
might be expressed in a very simple form by means of Sir W.\ Hamilton's theory of
quaternions. The property in question may be thus stated. Suppose a solid body
which revolves through an angle $\theta$ round an axis passing through the origin and
inclined to the axes of coordinates at angles $a$, $b$, $c$. Let
\[
\lambda = \tan\tfrac{1}{2}\theta\cos a, \qquad
\mu = \tan\tfrac{1}{2}\theta\cos b, \qquad
\nu = \tan\tfrac{1}{2}\theta\cos c,
\]
and write
\[
\Lambda = 1 + i\lambda + j\mu + k\nu;
\]
let $x$, $y$, $z$ be the coordinates of a point in the body previous to the rotation,
$x_{1}$, $y_{1}$, $z_{1}$, those of the same point after the rotation, and suppose
\begin{align*}
\Pi &= ix + jy + kz,\\
\Pi_{1} &= ix_{1} + jy_{1} + kz_{1};
\end{align*}
then the coordinates after the rotation may be determined by the formula
\[
\Pi_{1} = \Lambda\Pi\Lambda^{-1};
\]
viz., developing the second side of this equation,
\begin{align*}
\Pi_{1} = {}&i(\alpha\,x + \beta\,y + \gamma\,z)\\
&{}+ j(\alpha'x + \beta'y + \gamma'z)\\
&{}+ k(\alpha''x + \beta''y + \gamma''z),
\end{align*}

\vfill\pagebreak

% ===========================================================================
% PDF page 428 -- Book page 406
% ===========================================================================
where, putting to abbreviate $\kappa = 1 + \lambda^{2} + \mu^{2} + \nu^{2}$, we have
\begin{align*}
\kappa\alpha &= 1 + \lambda^{2} - \mu^{2} - \nu^{2}, & \kappa\alpha' &= 2(\lambda\mu + \nu), & \kappa\alpha'' &= 2(\lambda\nu - \mu),\\
\kappa\beta &= 2(\lambda\mu - \nu), & \kappa\beta' &= 1 - \lambda^{2} + \mu^{2} - \nu^{2}, & \kappa\beta'' &= 2(\mu\nu + \lambda),\\
\kappa\gamma &= 2(\lambda\nu + \mu), & \kappa\gamma' &= 2(\mu\nu - \lambda), & \kappa\gamma'' &= 1 - \lambda^{2} - \mu^{2} + \nu^{2};
\end{align*}
these values satisfying identically the well-known system of equations connecting the
quantities $\alpha$, $\beta$, $\gamma$, $\alpha'$, $\beta'$, $\gamma'$, $\alpha''$, $\beta''$, $\gamma''$.

\medskip

The quantities $a$, $b$, $c$, $\theta$ being immediately known when $\lambda$, $\mu$, $\nu$ are known, these
last quantities completely determine the direction and magnitude of the rotation, and
may therefore be termed the coordinates of the rotation; $\Lambda$ will be the quaternion of
the rotation. I propose here to develope a few of the consequences which may be
deduced from the preceding formulae.

\medskip

Suppose, in the first place, $\Pi = \Lambda - 1$, then $\Pi_{1} = \Lambda - 1$, which evidently implies
that the point is on the axis of rotation. The equation $\Pi_{1} = \Pi$ gives the identical
equations
\begin{align*}
\lambda(\alpha - 1) + \mu\beta + \nu\gamma &= 0,\\
\lambda\alpha' + \mu(\beta' - 1) + \nu\gamma' &= 0,\\
\lambda\alpha'' + \mu\beta'' + \nu(\gamma'' - 1) &= 0;
\end{align*}
from which, by changing the signs of $\lambda$, $\mu$, $\nu$, we derive
\begin{align*}
\lambda(\alpha - 1) + \mu\alpha' + \nu\alpha'' &= 0,\\
\lambda\beta + \mu(\beta' - 1) + \nu\beta'' &= 0,\\
\lambda\gamma + \mu\gamma' + \nu(\gamma'' - 1) &= 0.
\end{align*}
Hence evidently, whatever be the value of $\Pi$,
\[
\Lambda\Pi\Lambda^{-1} - \Pi = 0,
\]
if after the multiplication $i$, $j$, $k$ are changed into $\lambda$, $\mu$, $\nu$, a property which will be
required in the sequel.

\medskip

By changing the signs of $\lambda$, $\mu$, $\nu$, we also deduce
\begin{align*}
\Lambda^{-1}\Pi\Lambda = {}&i(\alpha x + \alpha'y + \alpha''z)\\
&{}+ j(\beta x + \beta'y + \beta''z)\\
&{}+ k(\gamma x + \gamma'y + \gamma''z),
\end{align*}
where $\alpha$, $\beta$, $\gamma$, $\alpha'$, $\beta'$, $\gamma'$, $\alpha''$, $\beta''$, $\gamma''$ are the same as before.

\medskip

Let the question be proposed to compound two rotations (both axes of rotation
being supposed to pass through the origin). Let $L$ be the first axis, $\Lambda$ the quaternion
of rotation, $L'$ the second axis, which is supposed to be fixed in space, so as not to
alter its direction by reason of the first rotation, $\Lambda'$ the corresponding quaternion of
rotation. The combined effect is given at once by
\[
\Pi_{1} = \Lambda'(\Lambda\Pi\Lambda^{-1})\Lambda'^{-1},
\]
that is,
\[
\Pi_{1} = \Lambda'\Lambda\Pi(\Lambda'\Lambda)^{-1};
\]

\vfill\pagebreak

% ===========================================================================
% PDF page 429 -- Book page 407
% ===========================================================================
or since (if $\Lambda_{1}$ be the quaternion for the combined rotation) $\Pi_{1} = \Lambda_{1}\Pi\Lambda_{1}^{-1}$, we have
clearly
\[
\Lambda_{1} = M_{1}\Lambda'\Lambda,
\]
$M_{1}$ denoting the reciprocal of the real part of $\Lambda'\Lambda$, so that
\[
M_{1}^{-1} = 1 - \lambda\lambda' - \mu\mu' - \nu\nu'.
\]
Retaining this value, the coefficients of the combined rotation are given by
\begin{align*}
\lambda_{1} &= M_{1}(\lambda + \lambda' + \mu'\nu - \mu\nu'),\\
\mu_{1} &= M_{1}(\mu + \mu' + \nu'\lambda - \nu\lambda'),\\
\nu_{1} &= M_{1}(\nu + \nu' + \lambda'\mu - \lambda\mu');
\end{align*}
to which may be joined [if $\kappa_{1} = 1 + \lambda_{1}^{2} + \mu_{1}^{2} + \nu_{1}^{2}$],
\[
\kappa_{1} = M_{1}^{2}\,\kappa'\kappa,
\]
$\kappa$, $\kappa'$, $\kappa_{1}$ as before. $\Lambda$ or $\Lambda'$ may be determined with equal facility in terms of
$\Lambda'$, $\Lambda_{1}$, or $\Lambda$, $\Lambda_{1}$. These formulae are given in my paper on the rotation of a solid
body (\textit{Cambridge Mathematical Journal}, vol.\ \textsc{iii}.\ p.\ 226, [6]).

\medskip

If the axis $L'$ be fixed in the body and moveable with it, its position after the
first rotation is obtained from the formula $\Pi_{1} = \Lambda\Pi\Lambda^{-1}$ by writing $\Pi = \Lambda' - 1$. Representing
by $\Lambda'' - 1$ the corresponding value of $\Pi_{1}$, we have $\Lambda'' = \Lambda\Lambda'\Lambda^{-1}$, which is the
value to be used instead of $\Lambda'$ in the preceding formula for the combined rotation,
thus the quaternion of rotation is proportional to $\Lambda\Lambda'\Lambda^{-1}\Lambda$, that is to $\Lambda\Lambda'$. Hence
here
\[
\Lambda_{1} = M_{1}\Lambda\Lambda',
\]
which only differs from the preceding in the order of the quaternion factors. If the
fixed and moveable axes be mixed together in any order whatever, the fixed axes
taken in order being $L$, $L'$, $\ldots$ and the moveable axes taken in order being $L_{0}$, $L_{0}'$ $\ldots$,
then the combined effect of the rotations is given by
\[
\Lambda_{1} = M \cdots \Lambda''\Lambda'\Lambda\Lambda_{0}\Lambda_{0}'\cdots,
\]
$M$ being the reciprocal of the real term of the product of all the quaternions.

\medskip

Suppose next the axes do not pass through the same point. If $\alpha$, $\beta$, $\gamma$ be the
coordinates of a point in $L$, and
\[
\Gamma = \alpha i + \beta j + \gamma k,
\]
then the formula for the rotation is
\[
\Pi_{1} - \Gamma = \Lambda(\Pi - \Gamma)\Lambda^{-1},
\]
or
\[
\Pi_{1} = \Lambda\Pi\Lambda^{-1} - (\Lambda\Gamma\Lambda^{-1} - \Gamma),
\]
where the first term indicates a rotation round a parallel axis through the origin, and
the second term a translation.

\vfill\pagebreak

% ===========================================================================
% PDF page 430 -- Book page 408
% ===========================================================================
For two axes $L$, $L'$ fixed in space,
\[
\Pi_{1} = \Lambda'\Lambda\Pi(\Lambda'\Lambda)^{-1} - (\Lambda'\Gamma'\Lambda'^{-1} - \Gamma') - \Lambda'(\Lambda\Gamma\Lambda^{-1} - \Gamma)\Lambda'^{-1};
\]
and so on for any number, the last terms being always a translation. If the two axes
are parallel, and the rotations equal and opposite,
\[
\Lambda = \Lambda'^{-1},
\]
whence
\[
\Pi_{1} = \Pi + \Lambda'(\Gamma - \Gamma')\Lambda'^{-1} - (\Gamma - \Gamma');
\]
or there is only a translation. The constant term vanishes if $i$, $j$, $k$ are changed into
$\lambda'$, $\mu'$, $\nu'$, which proves that the translation is in a plane perpendicular to the axes.

\medskip

Any motion of a solid body being represented by a rotation and a translation, it
may be required to resolve this into two rotations. We have
\[
\Pi_{1} = \Lambda_{1}\Pi\Lambda_{1}^{-1} + T,
\]
where $T$ is a given quaternion whose constant term vanishes. Hence, comparing this
with the general formula just given for the combination of two rotations,
\begin{align*}
\Lambda_{1} &= M_{1}\Lambda'\Lambda,\\
T &= -(\Lambda'\Gamma'\Lambda'^{-1} - \Gamma') - \Lambda'(\Lambda\Gamma\Lambda^{-1} - \Gamma)\Lambda'^{-1},
\end{align*}
the second of which equations may be simplified by putting $\Lambda'^{-1}T\Lambda' = S$, by which
it may be reduced to
\[
-S = (\Lambda'^{-1}\Gamma'\Lambda' - \Gamma') - (\Lambda\Gamma\Lambda^{-1} - \Gamma),
\]
which, with the preceding equation $\Lambda_{1} = M_{1}\Lambda'\Lambda$, contains the solution of the problem.
Thus if $\Lambda$ or $\Lambda'$ be given, the other is immediately known; hence also $S$ is known.
If in the last equation, after the multiplication is completely effected, we change
$i$, $j$, $k$ into $\lambda$, $\mu$, $\nu$, or $\lambda'$, $\mu'$, $\nu'$, we have respectively,
\[
S = \Lambda'^{-1}\Gamma'\Lambda' - \Gamma', \qquad S = -(\Lambda\Gamma\Lambda^{-1} - \Gamma),
\]
which are equations which must be satisfied by the coefficients of $\Gamma'$ and $\Gamma$ respectively.
Thus if the direction of one axis is given, that of the other is known, and the axes
must lie in certain known planes. If the position of one of the axes in its plane be
assumed, the equation containing $S$ divides itself into three others (equivalent to two
independent equations) for the determination of the position in its plane of the other
axis. If the axes are parallel, $\lambda$, $\mu$, $\nu$ are proportional to $\lambda'$, $\mu'$, $\nu'$; or changing
$i$, $j$, $k$ into $\lambda$, $\mu$, $\nu$, or $\lambda'$, $\mu'$, $\nu'$, we have $S = 0$; or what is the same thing, $T = 0$, which
shows that the translation must be perpendicular to the plane of the two axes.

\medskip

If $p$, $q$, $r$ have their ordinary signification in the theory of rotation, then from
the values in the paper in the \textit{Cambridge Mathematical Journal} already quoted,
\[
\kappa(ip + jq + kr) = 2\frac{d\Lambda}{dt}\Lambda + \frac{d\kappa}{dt};
\]

\vfill\pagebreak

% ===========================================================================
% PDF page 431 -- Book page 409
% ===========================================================================
but I have not ascertained whether this formula leads to any results of importance.
It may, however, be made use of to deduce the following property of quaternions, viz.,
if $\Lambda_{1} = M_{1}\Lambda'\Lambda$, $M_{1}$ as before, then
\[
\frac{1}{\kappa_{1}}\left(2\frac{d\Lambda_{1}}{dt}\Lambda_{1} + \frac{d\kappa_{1}}{dt}\right) = \frac{1}{\kappa}\left(2\frac{d\Lambda}{dt}\Lambda + \frac{d\kappa}{dt}\right),
\]
in which the coefficients of $\Lambda'$ are considered constant.

\medskip

To verify this \textit{a posteriori}, if in the first place we substitute for $\kappa_{1}$ its value
$M_{1}^{2}\kappa'\kappa$, we have
\[
\frac{d\kappa_{1}}{dt} = M_{1}^{2}\kappa'\frac{d\kappa}{dt} + 2\frac{\kappa_{1}}{M_{1}}\frac{dM_{1}}{dt}\kappa,
\]
and thence
\[
\frac{d\kappa_{1}}{dt}\Lambda_{1} = M_{1}^{3}\kappa'\frac{d\kappa}{dt}\Lambda'\Lambda + 2\frac{\kappa_{1}}{M_{1}}\frac{dM_{1}}{dt}\kappa\Lambda.
\]
Also
\[
\frac{d\Lambda_{1}}{dt}\Lambda_{1} = \left(\frac{1}{M_{1}}\frac{dM_{1}}{dt}\Lambda_{1} + M_{1}\Lambda'\frac{d\Lambda}{dt}\right)\Lambda_{1} = \frac{1}{M_{1}}\frac{dM_{1}}{dt}\Lambda_{1}^{2} + M_{1}^{2}\Lambda'\frac{d\Lambda}{dt}\Lambda'\Lambda,
\]
which reduces the equation to
\[
\Lambda_{1}^{2} + \kappa_{1} = 2\Lambda_{1} = 2M_{1}\Lambda'\Lambda,
\]
Hence observing that
\[
\Lambda_{1}^{2} + \kappa_{1} = 2\Lambda_{1} = 2M_{1}\Lambda'\Lambda,
\]
and omitting the factor $\Lambda$ from the resulting equation,
\[
\frac{2}{M_{1}^{2}}\frac{dM_{1}}{dt}\Lambda' + \Lambda'\frac{d\Lambda}{dt} = \kappa'\frac{d\Lambda}{dt};
\]
or since
\[
\frac{1}{M_{1}} = 1 - \lambda\lambda' - \mu\mu' - \nu\nu',
\]
substituting and dividing by $\Lambda'$, we obtain
\[
2\left(\lambda'\frac{d\lambda}{dt} + \mu'\frac{d\mu}{dt} + \nu'\frac{d\nu}{dt}\right)\Lambda'^{-1}\frac{d\Lambda}{dt} - \frac{d\Lambda}{dt}\Lambda'^{-1};
\]
or finally,
\begin{align*}
&2\left(\lambda'\frac{d\lambda}{dt} + \mu'\frac{d\mu}{dt} + \nu'\frac{d\nu}{dt}\right) = (1 - i\lambda' - j\mu' - k\nu')\left(i\frac{d\lambda}{dt} + j\frac{d\mu}{dt} + k\frac{d\nu}{dt}\right)\\
&\qquad - \left(i\frac{d\lambda}{dt} + j\frac{d\mu}{dt} + k\frac{d\nu}{dt}\right)(1 + i\lambda' + j\mu' + k\nu')\\
&\qquad = -(i\lambda' + j\mu' + k\nu')\left(i\frac{d\lambda}{dt} + j\frac{d\mu}{dt} + k\frac{d\nu}{dt}\right) - \left(i\frac{d\lambda}{dt} + j\frac{d\mu}{dt} + k\frac{d\nu}{dt}\right)(i\lambda' + j\mu' + k\nu'),
\end{align*}
which is obviously true.

\medskip
\hfill\textsc{c.}\hspace{1em}

\vfill\pagebreak

% ===========================================================================
% PDF page 432 -- Book page 410 -- Paper [69]: SUR LES DETERMINANTS GAUCHES
% ===========================================================================
\begin{center}
\textbf{\Large 69.}\\[1em]
\textbf{\large SUR LES D\'ETERMINANTS GAUCHES.}
\end{center}

\medskip

\begin{center}
\textit{[From the Journal f\"ur die reine und angewandte Mathematik (Crelle), tom.\ xxxviii.\ (1848),\\
pp.\ 93--96: continued from t.\ xxxii.\ p.\ 123, 52.]}
\end{center}

\medskip

J'ai nomm\'e ``\textit{D\'eterminant gauche}'' un d\'eterminant form\'e par un syst\`eme de
quantit\'es $\lambda_{r.s}$, qui satisfont aux conditions
\begin{equation}
\lambda_{r.s} = -\lambda_{s.r}\,(r + s). \tag{1}
\end{equation}
(o\`u les valeurs de $r$, $s$ s'\'etendent depuis l'unit\'e jusqu'\`a $n$).

\medskip

Or ces d\'eterminants peuvent facilement \^etre exprim\'es par un syst\`eme de d\'eterminants
pareils, dont les termes satisfont \`a ces conditions m\^eme dans le cas o\`u les
valeurs de $r$ et $s$ deviennent \'egales, ou pour lesquels on a
\begin{equation}
\lambda_{r.s} = -\lambda_{s.r}\,(r + s); \qquad \lambda_{r.r} = 0 \tag{2};
\end{equation}
ces d\'eterminants peuvent \^etre nomm\'es ``\textit{gauches et sym\'etriques}.''

\medskip

En effet, soit $\Omega$ le d\'eterminant gauche dont il s'agit, cette fonction peut \^etre
pr\'esent\'ee sous la forme
\begin{equation}
\Omega = \Omega_{0} + \Omega_{1}\lambda_{1.1} + \Omega_{2}\lambda_{2.2}\cdots + \Omega_{n}\lambda_{1.1}\lambda_{2.2}\cdots, \tag{3}
\end{equation}
o\`u $\Omega_{0}$ est ce que devient $\Omega$ si $\lambda_{1.1}$, $\lambda_{2.2}$, \&c.\ sont r\'eduits \`a z\'ero; $\Omega_{1}$ est ce que devient
le coefficient de $\lambda_{1.1}$, sous la m\^eme condition, et ainsi de suite; c'est-\`a-dire: $\Omega_{0}$ est
le d\'eterminant form\'e par les quantit\'es $\lambda_{r.s}$, en supposant que ces quantit\'es satisfassent
aux conditions (2), et en donnant \`a $r$ les valeurs $1, 2, 3, \ldots n$; $\Omega_{1}$ est le d\'eterminant
form\'e pareillement en donnant \`a $r$, $s$ les valeurs $2, 3, \ldots n$; $\Omega_{2}$ s'obtient en donnant
\`a $r$, $s$ les valeurs $1, 3, \ldots n$, et ainsi de suite; cela est ais\'e de voir si l'on range les
quantit\'es $\lambda_{r.s}$ en forme de carr\'e.

\vfill\pagebreak

% ===========================================================================
% PDF page 433 -- Book page 411
% ===========================================================================
Or les d\'eterminants gauches et sym\'etriques (savoir les d\'eterminants dont les termes
satisfont aux conditions (2)) se r\'eduisent \`a z\'ero pour un $n$ impair, et pour un $n$ pair
aux carr\'es des fonctions que M.\ Jacobi a trait\'ees dans son m\'emoire ``Ueber die Pfaff'-
sche Integrations-Methode'' (t.\ \textsc{ii}.\ [1827], p.\ 354, de ce journal) et dans le m\'emoire
``Theoria novi multiplicatoris aequationum differentialium'' t.\ \textsc{xxix}.\ [1845], p.\ 236, \&c.

\medskip

En effet, on voit d'abord (par ce que dit M.\ Jacobi) que le d\'eterminant s'\'evanouit
pour un $n$ impair, et que pour un $n$ pair il aura pour facteur la fonction dont il
s'agit; mais je ne sais pas si l'on a d\'ej\`a remarqu\'e que l'autre facteur se r\'eduit \`a la
m\^eme fonction.

\medskip

On obtient ces fonctions (dont je reprends ici la th\'eorie), par les propri\'et\'es g\'en\'erales
d'un d\'eterminant, d\'efini de la mani\`ere que voici: en exprimant par $(12 \ldots n)$ une
fonction quelconque dans laquelle entrent les nombres symboliques $1, 2, \ldots n$, et par $\pm$,
le signe correspondant \`a une permutation quelconque de ces nombres, la fonction
\[
\Sigma\pm(12 \ldots n)
\]
(o\`u $\Sigma$ d\'esigne la somme de tous les termes qu'on obtient en permutant ces nombres
d'une mani\`ere quelconque) est ce qu'on nomme D\'eterminant. On pourrait encore
g\'en\'eraliser cette d\'efinition en admettant plusieurs syst\`emes de nombres $1, 2, \ldots n$;
$1', 2', \ldots n'; \ldots$ qui alors devraient \^etre permut\'es ind\'ependamment les uns des autres;
on obtiendrait de cette mani\`ere une infinit\'e d'autres fonctions, mentionn\'ees (t.\ \textsc{xxx}.\ [1846]
p.\ 7). Dans le cas des d\'eterminants ordinaires, auquel je ne m'arr\^eterai pas ici, on aura
$(12 \ldots n) = \lambda_{1.1}\lambda_{2.2}\ldots\lambda_{n.n}$. Pour les cas des fonctions dont il s'agit (les fonctions de
M.\ Jacobi), on supposera $n$ pair, et l'on \'ecrira
\[
(12 \ldots n) = \lambda_{1.2}\lambda_{3.4}\ldots\lambda_{n-1.n},
\]
o\`u $\lambda_{r.s}$ sont des quantit\'es quelconques qui satisfont aux \'equations (1). La fonction sera
compos\'ee d'un nombre $1.2 \ldots n$ de termes; mais parmi eux il n'y aura que $1.3 \ldots (n - 1)$
termes diff\'erents qui se trouveront r\'ep\'et\'es $2^{\frac{1}{2}n}(1.2\ldots\tfrac{1}{2}n)$ fois, et qu'on obtiendra en
permutant cycliquement d'abord les $n - 1$ derniers nombres, puis les $n - 3$ derniers
nombres de chaque permutation, et ainsi de suite; le signe \'etant toujours $+$. Il
pourra \^etre d\'emontr\'e, comme pour les d\'eterminants, que ces fonctions changent de signe
en permutant deux quelconques des nombres symboliques, et qu'elles s'\'evanouissent si
deux de ces nombres deviennent identiques. De plus, en exprimant par $[12 \ldots n]$ la
fonction dont il s'agit, la r\`egle qui vient d'\^etre \'enonc\'ee, donnera pour la formation de
ces fonctions:
\[
[12 \ldots n] = \lambda_{1.2}[34 \ldots n] + \lambda_{1.3}[4 \ldots n-2]\ldots + \lambda_{1.n}[23 \ldots (n - 1)].
\]

\medskip

Cela pos\'e, revenons aux d\'eterminants gauches et sym\'etriques; et soit d'abord $n$ un
nombre impair. Alors le d\'eterminant sera compos\'e de plusieurs termes, chacun multipli\'e
par le produit de l'une des quantit\'es $\lambda_{1.2}$, $\lambda_{1.3}$, $\ldots \lambda_{1.n}$ par une des quantit\'es $\lambda_{2.1}$, $\lambda_{3.1}$, $\ldots \lambda_{n.1}$.
Il est facile de voir que pour chaque terme, multipli\'e par $\lambda_{1.\alpha}\lambda_{\beta.1}$ ($\alpha \neq \beta$), il existera
un terme \'egal et de signe contraire, multipli\'e par $\lambda_{1.\beta}\lambda_{\alpha.1}$. Or $\lambda_{1.\alpha}\lambda_{\beta.1} = \lambda_{\alpha.1}\lambda_{1.\beta}$: donc

\vfill\pagebreak

% ===========================================================================
% PDF page 434 -- Book page 412
% ===========================================================================
ces deux termes \textit{se d\'etruiront}. Restent les termes multipli\'es par $\lambda_{1.\alpha}\lambda_{\alpha.1}$: le coefficient
d'un terme de cette forme sera un d\'eterminant gauche de l'ordre $n - 2$, mais pr\'ecis\'ement
de la forme du d\'eterminant m\^eme de l'ordre $n$ dont il s'agit. Or $n$ \'etant impair,
$n - 2$ le sera \'egalement: donc tout d\'eterminant gauche et sym\'etrique s'\'evanouira, si
cela a lieu pour les d\'eterminants pareils d'un ordre inf\'erieur de deux unit\'es. Or pour
$n = 3$ on a \'evidemment $\lambda_{1.2}\lambda_{2.3}\lambda_{3.1} + \lambda_{1.3}\lambda_{2.1}\lambda_{3.2} = 0$, donc:

\medskip

``Tout d\'eterminant gauche et sym\'etrique d'un ordre impair est z\'ero.''

\medskip

Soit maintenant $n$ un nombre pair. Consid\'erons le d\'eterminant gauche et sym\'etrique
plus g\'en\'eral qu'on obtient en donnant \`a $r$ les valeurs $\alpha$, $2$, $3$, $\ldots n$, et \`a $s$ les
valeurs $\beta$, $2$, $3$, $\ldots n$. En d\'eveloppant cette fonction comme on vient de le faire dans
le cas d'un $n$ impair, il se pr\'esentera d'abord un terme multipli\'e par $\lambda_{\alpha.\beta}$. Mais ce
terme sera un d\'eterminant gauche et sym\'etrique de l'ordre $n - 1$ (savoir celui que l'on
obtiendrait en donnant \`a $r$, $s$ les valeurs $2, 3, \ldots n$), et comme $n - 1$ est impair, ce
terme s'\'evanouira. Puis le coefficient de $-\lambda_{\alpha.\alpha'}\lambda_{\beta'.\beta}$ (ou de $\lambda_{\alpha.\alpha'}\lambda_{\beta.\beta'}$) sera le d\'eterminant
gauche et sym\'etrique qu'on obtient en donnant \`a $r$ les valeurs $2, 3, \ldots n$ ($\alpha'$ except\'e),
et \`a $s$ les valeurs $2, 3, \ldots n$ ($\beta'$ except\'e). En admettant que ce d\'eterminant gauche
se r\'eduit \`a $[(\alpha'+1)\ldots n\,2\ldots(\alpha'-1)][(\beta'+1)\ldots n\,2\ldots(\beta'-1)]$, on aura
\[
\lambda_{\alpha.\alpha'}[(\alpha'+1)\ldots(\alpha'-1)]\cdot\lambda_{\beta.\beta'}[(\beta'+1)\ldots(\beta'-1)]
\]
pour le terme g\'en\'eral, et la somme de tous ces termes se r\'eduira \`a
\[
\{\lambda_{\alpha.2}[34\ldots n] + \lambda_{\alpha.3}[4\ldots n\,2] + \ldots\}\cdot\{\lambda_{\beta.2}[34\ldots n] + \lambda_{\beta.3}[4\ldots n\,2] + \ldots\},
\]
ou enfin \`a $[\alpha 23\ldots n][\beta 23\ldots n]$. Or si le th\'eor\`eme en question a lieu pour $n - 2$,
il aura lieu aussi pour $n$. Pour $n = 2$ le d\'eterminant est $\lambda_{\alpha.\beta}\lambda_{2.2} - \lambda_{\alpha.2}\lambda_{\beta.2}$, c'est-\`a-dire
$\lambda_{\alpha.2}\lambda_{\beta.2}$ ou $[\alpha 2]\cdot[\beta 2]$: donc la m\^eme chose a toujours lieu et on obtient le th\'eor\`eme
que voici:

\medskip

``Le d\'eterminant gauche et sym\'etrique qu'on obtient en donnant \`a $r$ les valeurs
$\alpha$, $2$, $3$, $\ldots n$, et \`a $s$ les valeurs $\beta$, $2$, $3$, $\ldots n$ (o\`u $n$ est pair), se r\'eduit \`a
\[
[\alpha 23\ldots n]\cdot[\beta 23\ldots n];
\]
et en particulier, en donnant \`a $r$, $s$ les valeurs $1, 2, \ldots n$, ce d\'eterminant se r\'eduit \`a
\[
[12\ldots n]^{2}.\text{''}
\]

\medskip

Appliquons ces th\'eor\`emes \`a la r\'eduction de l'\'equation (3). En supposant pour
plus de simplicit\'e $\lambda_{r.r} = 1$, on aura pour un $n$ pair:
\[
\Omega = [123\ldots n]^{2} + [34\ldots n]^{2} + [56\ldots n]^{2} + \ldots + 1;
\]
\[
+ [24\ldots n]^{2}
\]

\vfill\pagebreak

% ===========================================================================
% PDF page 435 -- Book page 413
% ===========================================================================
et pour un $n$ impair:
\[
\Omega = [23\ldots n]^{2} + [45\ldots n]^{2} + \ldots + 1.
\]
\[
+ [13\ldots n]^{2}
\]

\medskip

Particuli\`erement pour $n = 4$ on obtient:
\begin{align*}
\Omega &= [1234]^{2} + [12]^{2} + [13]^{2} + [14]^{2} + [23]^{2} + [34]^{2} + [42]^{2} + 1,\\
&= (\lambda_{1.2}\lambda_{3.4} + \lambda_{1.3}\lambda_{4.2} + \lambda_{1.4}\lambda_{2.3})^{2} + \lambda_{1.2}^{2} + \lambda_{1.3}^{2} + \lambda_{1.4}^{2} + \lambda_{3.4}^{2} + \lambda_{4.2}^{2} + \lambda_{2.3}^{2} + 1;
\end{align*}
et pour $n = 3$:
\begin{align*}
\Omega &= [23]^{2} + [31]^{2} + [12]^{2} + 1,\\
&= \lambda_{2.3}^{2} + \lambda_{3.1}^{2} + \lambda_{1.2}^{2} + 1;
\end{align*}
r\'esultats qui s'accordent parfaitement avec ceux que j'ai donn\'es dans mon premier
m\'emoire sur ce sujet.

\medskip

Il serait possible de trouver pour les quantit\'es $\Lambda_{r.s}$ (m\'emoire cit\'e) des expressions
analogues \`a celles que nous venons de donner pour $\Omega$: mais elles seraient beaucoup plus
compliqu\'ees.

\vfill\pagebreak

% ===========================================================================
% PDF page 436 -- Book page 414 -- Paper [70] begins
% ===========================================================================
\begin{center}
\textbf{\Large 70.}\\[1em]
\textbf{\large SUR QUELQUES TH\'EOR\`EMES DE LA G\'EOM\'ETRIE DE POSITION.}
\end{center}

\medskip

\begin{center}
\textit{[From the Journal f\"ur die reine und angewandte Mathematik (Crelle), tom.\ xxxviii.\ (1848),\\
pp.\ 97--104: continued from t.\ xxxiv.\ p.\ 275, 55.]}
\end{center}

\medskip

\begin{center}\S\ I.\end{center}

\medskip

Le th\'eor\`eme de Pascal appliqu\'e au cas o\`u une conique se r\'eduit \`a deux droites,
donne lieu \`a un syst\`eme de neuf points situ\'es trois \`a trois dans neuf droites qui
passent r\'eciproquement trois \`a trois par les neuf points. Je repr\'esenterai ces points par
\[
1, 4, 7;\quad 2, 5, 8;\quad 3, 6, 9,
\]
et les droites par
\[
135, 426, 789;\quad 129, 483, 756;\quad 186, 459, 723.
\]

\medskip

On peut prendre pour la conique dont il s'agit, deux quelconques de ces droites
qui ne se rencontrent pas dans un des neuf points, ou ce qui revient au m\^eme, deux
quelconques des droites qui appartiennent \`a un des trois syst\`emes dans lesquels les
droites viennent d'\^etre divis\'ees; alors la troisi\`eme droite sera celle qui contient les points
de rencontre des c\^ot\'es oppos\'es de l'hexagone. Par exemple, en prenant pour conique
les deux droites $135$, $246$, l'hexagone sera $123456$, et les points de rencontre des c\^ot\'es
oppos\'es, savoir de $12$ et $45$, de $23$ et $56$, et de $34$ et $61$, seront respectivement $9$, $7$, $8$;
c'est-\`a-dire des points situ\'es dans la droite $789$; de mani\`ere que:

\medskip

TH\'EOR\`EME XV. ``La figure form\'ee par neuf points situ\'es 3 \`a 3 dans 9 droites peut
\^etre consid\'er\'ee, de neuf mani\`eres diff\'erentes, comme r\'esultante du th\'eor\`eme de Pascal
appliqu\'e au cas o\`u la conique se r\'eduit \`a deux droites.''

\medskip

Il y a, comme l'a remarqu\'e M.\ Graves dans le m\'emoire cit\'e du \textit{Philosophical
Magazine}, une autre mani\`ere assez singuli\`ere d'envisager la figure. Consid\'erons pour
cela les trois triangles $126$, $489$, $735$, qui peuvent \^etre d\'eriv\'es assez simplement de

\vfill\pagebreak

% ===========================================================================
% PDF page 437 -- Book page 415
% ===========================================================================
l'arrangement $135$, $426$, $789$. Le premier de ces triangles est circonscrit au second; car
les c\^ot\'es $26$, $61$, $12$ contiennent respectivement les points $4$, $8$, $9$; \'egalement le second
est circonscrit au troisi\`eme; et le troisi\`eme au premier. On tire encore du m\^eme
arrangement $135$, $426$, $789$ un autre syst\`eme pareil de triangles; savoir $189$, $435$, $726$.
Pareillement les arrangements $129$, $483$, $756$ et $186$, $459$, $723$ donnent lieu chacun \`a
deux syst\`emes analogues. Donc:

\medskip

TH\'EOR\`EME XVI. La figure form\'ee par neuf points situ\'es 3 \`a 3 dans 9 droites
peut \^etre consid\'er\'ee, de six mani\`eres diff\'erentes, comme compos\'ee de trois triangles
circonscrits cycliquement l'un \`a l'autre.

\medskip

Repr\'esentons maintenant par $1, 2, 3, 4, 5, 6, 7, 8, 9$ les neuf points d'inflexion
d'une courbe du troisi\`eme ordre (six de ces points seront n\'ecessairement imaginaires).
Ces points sont, comme on sait, situ\'es 3 \`a 3 dans douze droites qui passent r\'eciproquement
quatre \`a quatre par les neuf points, et qui peuvent \^etre repr\'esent\'ees par
\[
135, 426, 789;\quad 129, 483, 756;\quad 186, 459, 723;\quad 147, 258, 369;
\]
de sorte que ce syst\`eme est un cas particulier de celui que nous venons de consid\'erer.
En consid\'erant deux quelconques des droites qui ont un point en commun, par exemple
$147$, $186$, et les deux autres droites qui passent par ce m\^eme point $135$, $129$, on obtient
par l\`a deux quadrilat\`eres $4876$ et $3952$ qui ont pour centre commun le point $7$ et
dont chacun est circonscrit \`a l'autre. On aurait pu obtenir par ces m\^emes droites
$147$, $186$, $135$, $129$ des autres syst\`emes pareils; donc

\medskip

TH\'EOR\`EME XVII. Huit points quelconques, parmi neuf points d'inflexion d'une
courbe du troisi\`eme ordre, peuvent \^etre consid\'er\'es de trois mani\`eres diff\'erentes comme
formant deux quadrilat\`eres circonscrits l'un \`a l'autre. Les diagonales de ces six
quadrilat\`eres se r\'eduisent \`a quatre droites qui passent par le neuvi\`eme point d'inflexion.

\medskip

\begin{center}\S\ VI.\quad \textit{Sur les figures r\'eciproques.}\end{center}

\medskip

Soient $\xi : \omega$, $\eta : \omega$, $\rho : \omega$ les coordonn\'ees d'un point. Les deux plans d\'efinis par
les \'equations
\[
\begin{aligned}
&(a\xi + b\eta + c\rho + d\omega)x\\
&+ (a'\xi + b'\eta + c'\rho + d'\omega)y\\
&+ (a''\xi + b''\eta + c''\rho + d''\omega)z\\
&+ (a'''\xi + b'''\eta + c'''\rho + d'''\omega)w
\end{aligned} = 0, \qquad
\text{et} \qquad
\begin{aligned}
&(a\xi + a'\eta + a''\rho + a'''\omega)x\\
&+ (b\xi + b'\eta + b''\rho + b'''\omega)y\\
&+ (c\xi + c'\eta + c''\rho + c'''\omega)z\\
&+ (d\xi + d'\eta + d''\rho + d'''\omega)w
\end{aligned} = 0,
\]
seront les plans r\'eciproques du point donn\'e.

\medskip

Il peut arriver que le plan r\'eciproque d'un point quelconque passe par ce point
m\^eme; ce qui implique aussi l'identit\'e des deux plans r\'eciproques. Ce cas particulier
a \'et\'e l'objet des recherches de M.\ M\"obius. Je parlerai dans la suite des r\'eciproques
de cette esp\`ece en les appelant \textit{R\'eciproques gauches}.

\medskip

Mais g\'en\'eralement, pour que les plans r\'eciproques d'un point passent par ce point
m\^eme, il faut que le point soit situ\'e sur une certaine surface du second ordre; et

\end{document}
